\chapter{Introdução}\label{cap:introducao}

As doenças hepáticas representam um grave problema de saúde pública mundial, sendo responsáveis por milhões de mortes anuais devido a complicações como cirrose, hepatite viral e carcinoma hepatocelular \citep{sousa2023}. No Brasil, estudos recentes indicam que a mortalidade por hepatites crônicas aumentou nas regiões Norte e Nordeste, enquanto a mortalidade por hepatites A, B, outras virais e não especificadas apresentou tendência de queda entre 2001 e 2020 \citep{sousa2023}. Esses dados evidenciam a necessidade de monitoramento longitudinal da evolução clínica dos pacientes.

O diagnóstico e acompanhamento de pacientes hepatopatas dependem fortemente de registros clínicos detalhados, armazenados em prontuários eletrônicos. A digitalização da assistência e a adoção ampla desses prontuários produziram grandes volumes de dados longitudinais, que documentam a trajetória clínica de cada paciente ao longo de meses ou anos. Grande parte desse acervo, contudo, não é estruturada: estima-se que cerca de 80\% dos dados médicos permaneçam não estruturados e subaproveitados após sua criação \citep{kong2019}, dispersos em linguagem natural, o que dificulta sua análise sistemática \citep{schneider2023}. A falta de ferramentas computacionais eficientes impacta diretamente a qualidade do acompanhamento clínico, prejudicando a detecção precoce da progressão da doença e o planejamento de estratégias de tratamento \citep{hou2014}.

O Processamento de Linguagem Natural (PLN) tem sido amplamente explorado para extração de informações médicas. Revisões recentes apontam avanços significativos no uso de PLN em gastroenterologia e hepatologia, incluindo técnicas baseadas em grandes modelos de linguagem (LLMs), mas ainda existem desafios relacionados a abreviações, terminologia especializada e evolução temporal dos prontuários \citep{stammers2025,shahab2024}. Ferramentas como ChatGPT têm mostrado potencial para auxiliar na síntese de informações médicas, embora com limitações quanto a precisão e vieses \citep{dave2023}, e LLMs conversacionais passaram a ser avaliados como ferramentas automatizadas de apoio à decisão clínica, como na interpretação de biomarcadores complexos em oncologia de precisão, tarefa que hoje exige a consulta manual a estudos e bases de dados prévios \citep{benary2023}.

Nos últimos anos, técnicas de Retrieval-Augmented Generation (RAG)\citep{gao2024} têm sido aplicadas para integrar extração de conhecimento e recuperação de informações em bases especializadas. Entre essas, o GraphRAG combina LLMs com estruturas de grafos, permitindo modelar hierarquias e relações complexas entre entidades clínicas, o que facilita a análise ao longo do tempo e a geração de resumos contextualizados \citep{edge2024,xiang2025}. Estudos recentes destacam que a eficácia do GraphRAG depende do tipo de tarefa, da qualidade do grafo e da estratégia de construção de comunidades, sendo crucial avaliar seu desempenho em cenários clínicos \citep{xiang2025}. Apesar dos avanços no uso de técnicas de  PLN para a extração de informações médicas em prontuários eletrônicos, a análise longitudinal da evolução clínica de pacientes hepatopatas ainda apresenta desafios significativos.

\section{Objetivo Geral}

 Nesse contexto, o principal objetivo deste trabalho é desenvolver e avaliar uma abordagem baseada em GraphRAG para a extração e o monitoramento de informações clínicas em prontuários eletrônicos de pacientes hepatopatas, buscando responder aos seguintes questionamentos:

\begin{enumerate}
    \item Como a aplicação da sumarização de contexto pode melhorar a análise de informações clínicas relevantes em prontuários eletrônicos de pacientes hepatopatas?
    \item Em que medida o modo de busca do GraphRAG e o nível de granularidade da detecção de comunidades influenciam a identificação correta das doenças hepáticas documentadas no prontuário?
    \item Quais são as limitações e desafios na integração de GraphRAG com dados clínicos não estruturados em termos de acurácia e relevância das informações extraídas?
\end{enumerate}

\subsection{Objetivos Específicos}

Para atingir o objetivo geral, este trabalho se propõe a:

\begin{enumerate}
    \item Construir um conjunto de prontuários rotulado para as oito doenças hepáticas de interesse, por meio de um funil de triagem textual, triagem por LLM e anotação manual;

    \item Indexar o corpus com o GraphRAG, com ajuste do \textit{prompt} de extração de entidades e relações ao domínio da hepatologia, e construir a hierarquia de comunidades pelo algoritmo de Leiden;

    \item Desenvolver e refinar, ao longo de versões sucessivas, o \textit{prompt} de consulta que extrai da resposta do modelo as doenças hepáticas documentadas em cada prontuário;

    \item Avaliar sistematicamente os quatro modos de busca do GraphRAG (\textit{basic}, \textit{local}, \textit{global} e \textit{drift}) em dois níveis de comunidade, por meio de métricas de classificação multirrótulo e testes estatísticos, contra o \textit{ground truth} anotado;

    \item Investigar, por meio de testes estatísticos e análise qualitativa, os mecanismos que explicam os casos de acerto e de erro de cada modo de busca, em particular do modo \textit{drift}.

\end{enumerate}

\section{Relevância do Trabalho}\label{sec:relevancia}

As doenças hepáticas representam uma significativa carga global de saúde, sendo responsáveis por elevado número de internações, complicações graves e mortalidade, a literatura aponta para a necessidade de transição da prática tradicional para a medicina de precisão em hepatologia, graças ao volume crescente de dados clínicos, laboratoriais, de imagem e genômicos disponíveis \citep{SABRIE2023350}. No Brasil, a mortalidade por hepatites crônicas aumentou em algumas regiões (Norte e Nordeste) enquanto outras formas de hepatite apresentaram tendência de queda entre 2001 e 2020, reforçando a urgência de ferramentas de monitoramento de evolução clínica \citep{sousa2023}.

Em paralelo, os prontuários eletrônicos e notas clínico-assistenciais contêm riqueza de dados em texto livre, que até o momento permanecem subutilizados devido à complexidade do processamento linguístico, revisões recentes em hepatologia e gastroenterologia ressaltam que o PLN já consegue extrair informações úteis desses textos, mas ainda padece de limitações como falta de generalização, baixa explicabilidade, escassez de integração temporal e escopo restrito \citep{omar2025}.

Nesse cenário, a combinação de PLN com grafos de conhecimento e modelos de recuperação assistida, como a abordagem GraphRAG, surge como uma oportunidade inovadora para capturar relações complexas entre entidades clínicas e permitir consultas mais ricas e estruturadas, algo ainda raro no contexto de hepatologia \citep{stammers2025}. Portanto, este trabalho se mostra relevante ao propor uma solução que aborda duas lacunas principais: (1) a extração inteligente de dados textuais não estruturados em hepatopatias, e (2) o uso de grafos de conhecimento para dar sentido e facilitar a exploração dessas informações por clínicos, pesquisadores e gestores. Essa relevância é tanto prática, ao contribuir para a melhoria no acompanhamento, prognóstico e tratamento de pacientes com doenças hepáticas, quanto metodológica, ao aplicar técnicas de ponta em PLN e grafos de conhecimento num domínio ainda pouco explorado.

% -----------------------------
% Bibliografia
% -----------------------------

